\documentclass[
    language=english,
    doctype=art-catalogue,
    institution=none,
]{../../omnilatex}

\RequirePackage{config/document-settings}

\begin{document}
\maketitle

\vspace{1cm}

\begin{center}
\textbf{BETWEEN WORLDS}\\[0.3cm]
\textit{The Art of Elena Vasquez}\\[0.5cm]
\textbf{Curated by Dr. Margaret Liu}\\[0.3cm]
Whitfield Gallery of Contemporary Art\\
June 15 --- August 30, 2024
\end{center}

\vspace{1cm}

\section*{FOREWORD}

Elena Vasquez makes art that lives in the spaces between things---between memory and reality, between the natural and the artificial, between the world we see and the world we imagine. Her work asks questions that have no easy answers, and it does so with a beauty that is both disarming and unsettling.

This exhibition, the largest retrospective of Vasquez's work to date, brings together forty-three pieces spanning twenty years of artistic practice. From the early watercolors that first brought her to public attention to the immersive installations that have made her one of the most important artists of her generation, \textit{Between Worlds} traces the arc of an extraordinary career.

The works in this exhibition are arranged not chronologically but thematically, organized around the central preoccupation of Vasquez's art: the boundary between the seen and the unseen, the known and the unknown, the real and the imagined.

--- Dr. Margaret Liu, Curator

\section*{ARTIST BIOGRAPHY}

Elena Vasquez was born in Santa Fe, New Mexico, in 1974. She studied fine art at the University of New Mexico and the Royal College of Art in London, where she developed her distinctive approach to mixed media and installation art.

Vasquez's early work was rooted in traditional watercolor and printmaking, but she quickly moved beyond conventional mediums, incorporating found objects, natural materials, and digital technology into her practice. Her work has been exhibited at institutions including the Museum of Modern Art (New York), the Tate Modern (London), the Centre Pompidou (Paris), and the Guggenheim Museum (New York).

She has received numerous awards, including the MacArthur Fellowship (2012), the Hugo Boss Prize (2015), and the Praemium Imperiale (2019). Her work is held in the permanent collections of major museums worldwide.

Vasquez currently lives and works in Marfa, Texas, where she maintains a studio in a converted military barracks on the edge of the Chihuahuan Desert.

\section*{EXHIBITION CATALOGUE}

\subsection{Gallery 1: Thresholds}

\textbf{Plate 1:} \textit{The Door That Wasn't There Before} (2019)
\begin{quote}
Mixed media installation. Wooden door frame, cast resin, LED lighting, sound system.
$180 \times 90 \times 45$ cm.

A freestanding door frame stands in the center of the gallery, its surface covered in layers of cast resin that create the illusion of depth---as though the door opens not onto another room but onto another dimension. Embedded speakers emit a low, continuous hum that changes pitch as viewers approach, creating the sensation that the door is listening.

The work was inspired by a recurring dream Vasquez had as a child, in which she would wake to find a new door in her bedroom that led to a world made entirely of light. ``I spent my childhood looking for that door,'' she has said. ``This work is my attempt to build it.''
\end{quote}

\textbf{Plate 2:} \textit{Map of a Country That Doesn't Exist} (2020)
\begin{quote}
Ink and gold leaf on handmade paper.
$120 \times 90$ cm.

A large-scale map rendered in exquisite detail, depicting a country with cities, rivers, mountain ranges, and political boundaries---none of which correspond to any real place. The map includes a legend, a scale, and a compass rose, all executed with the precision of actual cartography.

Vasquez spent three years constructing this map, developing an entire fictional geography complete with its own history, culture, and political system. ``Every country is an invention,'' she has said. ``Some are just more honest about it than others.''
\end{quote}

\textbf{Plate 3:} \textit{The Weight of Light} (2018)
\begin{quote}
Cast concrete, gold leaf, fiber optic cable.
$60 \times 60 \times 60$ cm.

A solid cube of concrete, its surface embedded with fiber optic cables that emit a soft, warm glow. The top surface is covered in gold leaf, creating a pool of reflected light that shifts and changes as viewers move around the piece.

The work explores the paradox of light as both weightless and transformative. ``Light has no mass,'' Vasquez notes, ``but it changes everything it touches. That seems like a kind of weight to me.''
\end{quote}

\subsection{Gallery 2: Erosion}

\textbf{Plate 4:} \textit{What the River Remembers} (2021)
\begin{quote}
Video installation, 3 channels. Duration: 47 minutes (looping).
Three screens arranged in a triangular formation, each displaying footage of the same river filmed from different perspectives---above, below, and at water level. The footage is time-lapsed, showing the river's seasonal changes compressed into a single cycle: ice, flood, drought, ice again.

The sound design layers recordings of the river at different points along its course, creating a sonic landscape that is both familiar and alien. The work runs continuously, with no beginning or end, inviting viewers to enter at any point in the cycle.
\end{quote}

\textbf{Plate 5:} \textit{Erosion Study \#7} (2020)
\begin{quote}
Sandstone, water, steel.
Variable dimensions.

A large slab of sandstone sits in a shallow steel basin. A thin stream of water flows continuously over the stone's surface, slowly wearing it away over the course of the exhibition. By the closing date, the stone will have lost approximately 2 centimeters of material, its surface transformed by the patient, relentless action of water.

The work is a meditation on time and impermanence. ``Everything is eroding,'' Vasquez observes. ``The question is what we choose to preserve, and what we allow to change.''
\end{quote}

\textbf{Plate 6:} \textit{Ghost Forest} (2022)
\begin{quote}
Reclaimed wood, steel wire, LED lighting.
$400 \times 300 \times 300$ cm (installation).

Twelve reclaimed wooden posts, each stripped of bark and polished smooth, arranged in a circular formation. Steel wires connect the posts at various heights, creating a web-like structure that catches the light from embedded LEDs, casting intricate shadows on the gallery walls.

The posts were salvaged from a demolished warehouse in Detroit. ``These posts held up a building for fifty years,'' Vasquez explains. ``Now they hold up nothing but light and air. I think there's something beautiful in that---in becoming something else after you've stopped being what you were meant to be.''
\end{quote}

\subsection{Gallery 3: Thresholds of Perception}

\textbf{Plate 7:} \textit{The Color of Silence} (2023)
\begin{quote}
Sound-absorbing panels, pigment, steel frame.
$250 \times 250$ cm (wall-mounted).

A large, square panel covered in sound-absorbing material, its surface treated with a single pigment that shifts from deep blue at the center to pale white at the edges, creating the illusion of infinite depth. The panel absorbs all sound in its vicinity, creating a zone of absolute silence that viewers can step into and experience.

The work challenges the primacy of vision in art. ``We assume that art is something we look at,'' Vasquez says. ``But what if it's something we hear---or don't hear? Silence is the most overlooked medium there is.''
\end{quote}

\textbf{Plate 8:} \textit{Blind Landscape} (2022)
\begin{quote}
Clay, beeswax, essential oils.
Dimensions variable.

A series of twelve clay tablets, each embedded with a different blend of essential oils and coated in beeswax. The tablets are displayed in a darkened room, and viewers are invited to touch them and experience the landscape through scent and texture rather than sight.

The work was developed during a residency in Provence, where Vasquez spent months walking through lavender fields and olive groves with her eyes closed, learning to navigate by scent and sound. ``I wanted to create a landscape you could feel,'' she explains. ``One that exists in the body, not just in the eye.''
\end{quote}

\textbf{Plate 9:} \textit{Self-Portrait as a Room} (2024)
\begin{quote}
Wood, mirrors, sound recording, fabric.
$300 \times 400 \times 250$ cm (installation).

A small wooden structure, just large enough for one person to enter. The interior is lined with mirrors, and a hidden speaker plays a continuous recording of the artist breathing. Viewers who enter the structure see themselves reflected infinitely, surrounded by the sound of another person's breath---a disorienting experience that blurs the boundary between self and other.

``This is the most personal work I've ever made,'' Vasquez says. ``It's a portrait of what it feels like to be inside my own head---surrounded by reflections, haunted by the sound of my own breathing.''
\end{quote}

\section*{SELECTED BIBLIOGRAPHY}

\begin{enumerate}
    \item Liu, Margaret. \textit{Elena Vasquez: Between Worlds}. Whitfield Gallery of Contemporary Art, 2024.
    \item Collins, Judith. \textit{The Art of Impermanence: Elena Vasquez and the Poetics of Erosion}. Thames \& Hudson, 2022.
    \item Okafor, Chinua. ``The Weight of Light: An Interview with Elena Vasquez.'' \textit{Artforum}, vol. 61, no. 4, 2023, pp. 112--121.
    \item Vasquez, Elena. \textit{Map of a Country That Doesn't exist: Selected Writings}. Siglio Press, 2021.
    \item Park, Soojin. ``Ghost Forest: Reclamation and Transformation in the Work of Elena Vasquez.'' \textit{October}, no. 189, 2022, pp. 45--67.
    \item Rivera, Carlos. ``The Color of Silence: Elena Vasquez at the Venice Biennale.'' \textit{Frieze}, no. 238, 2023, pp. 88--93.
\end{enumerate}

\section*{EXHIBITION DETAILS}

\begin{table}[h]
\centering
\begin{tabular}{lr}
\toprule
\textbf{Venue} & Whitfield Gallery of Contemporary Art \\
\textbf{Dates} & June 15 --- August 30, 2024 \\
\textbf{Hours} & Tuesday--Sunday, 10:00 AM--6:00 PM \\
\textbf{Admission} & Free \\
\textbf{Curator} & Dr. Margaret Liu \\
\textbf{Exhibition Design} & Studio vasquez \\
\textbf{Catalogue} & ISBN 978-0-12345-678-9 \\
\bottomrule
\end{tabular}
\caption{Exhibition information}
\end{table}

\vspace{1cm}

\begin{center}
\textit{This exhibition is made possible by the generous support of\\
the National Endowment for the Arts, the Andy Warhol Foundation,\\
and the Whitfield Gallery Patrons Circle.}
\end{center}

\end{document}
